\section{Validation}
\label{sec:validation}

\subsection{Against \taxsim{}}
\label{sec:validation-taxsim}

\policyengineus{} is cross-validated against NBER's \taxsim{}35 on Current
Population Survey tax units. Two independent artifacts carry this
comparison:

\begin{enumerate}
\item \textbf{policyengine-us's own validation notebook}
(\texttt{docs/validation/taxsim.ipynb}) computes, for the 2022 \cps{}, the
share of tax units whose \policyengineus{}-computed federal income tax (and
state income tax, for Massachusetts, New York, and Washington) falls within
\$100 of \taxsim{}35's. \todo{Regenerate this notebook's output at
manuscript-drafting time and cite the current numbers -- the version seen
during outline drafting (2026-07-04) is not asserted here because it is not
this paper's committed evidence artifact; running the notebook is a
render-only step (no new modeling), so it belongs in v1 rather than as a
TODO(evidence needed).}
\evidence{/Users/maxghenis/PolicyEngine/policyengine-us/docs/validation/taxsim.ipynb}
\item \textbf{The policyengine-taxsim emulator repository} maintains
checked-in, per-household comparison tables against \taxsim{} for tax years
2021--2024 (6{,}000 households per year), with exact-match boolean columns
for federal, state, and combined liability
(\texttt{federal\_match}, \texttt{state\_match}, \texttt{overall\_match}).
These are strict exact-match indicators, not tolerance bands, and should be
reported alongside the \$100-tolerance figures from item 1 with that
distinction stated explicitly rather than pooled into one number.
\evidence{/Users/maxghenis/PolicyEngine/policyengine-taxsim/dashboard/build/data/\{2021,2022,2023,2024\}/comparison\_results\_\{year\}.csv}
\end{enumerate}

\texttt{programs.yaml} records \taxsim{} validation as a property of two
specific programs -- federal income tax (``Validated against NBER TAXSIM for
2022'') and state income tax (``Validated against NBER TAXSIM for
2021+'') -- which this section's numbers should extend rather than
duplicate.
\evidence{/Users/maxghenis/PolicyEngine/policyengine-us/policyengine\_us/programs.yaml}

\todo{Open methodological question for v1: the two artifacts use different
comparison conventions (\$100 tolerance vs.\ exact match) and appear to have
been generated independently rather than from one shared script. Before
submission, either reconcile them into one table with both conventions
reported side by side, or explain why they measure different things (e.g.
one runs on 2022 only, the other spans 2021--2024; check whether they use
the same \taxsim{} version and the same CPS extract before treating them as
comparable).}

\subsection{Against administrative aggregates}
\label{sec:validation-admin}

\populace{}'s US build includes an out-of-sample backtest of baseline levels
against published administrative totals for tax concepts that are
\emph{not} calibration targets: nine benchmark rows compare
\policyengineus{}'s baseline weighted total for a variable (at the dataset's
2024 period) against IRS \soi{} Individual Complete Report (Publication
1304) Tax Year 2023 actuals, from Tables 3.3 and 1.4, with an explicitly
documented one-year vintage gap.
\evidence{/Users/maxghenis/PolicyEngine/populace/packages/populace-build/src/populace/build/us/soi\_baseline\_levels.json;
example benchmark row \texttt{soi\_income\_tax\_net}: benchmark value
\$2{,}042{,}047{,}899{,}000 (TY2023), source IRS SOI Pub 1304 Table 3.3,
url \url{https://www.irs.gov/pub/irs-soi/23in33ar.xls}}

\todo{Evidence needed: the \emph{result} of this backtest (populace's actual
baseline total versus each of the nine benchmarks) is not committed in the
working tree as of outline drafting -- \texttt{soi\_baseline\_levels.json}
defines the benchmark targets and comparison methodology, but running the
comparison against a specific certified release is new (light) analysis,
not assembly. Run it against the release pinned in
Section~\ref{sec:disclosures} and report all nine rows, not a favorable
subset.}

\subsection{Against \jct{} and \cbo{} published scores}
\label{sec:validation-jct-cbo}

The model's own methodology documentation currently lists validation
against \taxsim{} and IRS \soi{} as a planned (not yet written) section,
with no mention of \jct{} or \cbo{}.
\evidence{/Users/maxghenis/PolicyEngine/policyengine-us/docs-quarto/methodology/index.qmd:
"Planned sections (not yet written): microdata construction (enhanced CPS),
calibration to administrative totals, behavioral responses, uprating, and
validation against NBER TAXSIM and IRS SOI."} Existing \populace{} build
artifacts already go further than that documentation reflects, tying
specific \policyengineus{}-modeled provisions to named, dated, published
\jct{} documents:

\begin{itemize}
\item \textbf{OBBBA provisions vs.\ JCX-35-25.} Nineteen provisions of the
2025 One Big Beautiful Bill Act that \policyengineus{} models in its
baseline are each encoded as a counterfactual patch (turning that provision
off against a pre-OBBBA / TCJA-expiration reference policy), paired with
\jct{}'s own FY2026 conventional revenue score for that provision from
JCX-35-25 (published 2025-07-01).
\evidence{/Users/maxghenis/PolicyEngine/populace/packages/populace-build/src/populace/build/us/obbba\_reforms.json;
\texttt{jct\_document}: \{id: "JCX-35-25", published: "2025-07-01", url:
jct.gov/publications/2025/jcx-35-25/\}; example row
\texttt{obbba\_reduced\_rates}: JCT FY2026 score -\$147{,}679{,}000{,}000}
The scoring engine that runs these 19 provisions through the
\policyengineus{} microsimulation and emits a comparison
(\texttt{reform\_validation.json}, labeled in-sample vs.\ out-of-sample) is
committed, but \emph{its output for a specific release is not} -- it is
generated at release-build time, published alongside the release bundle
rather than checked into the populace working tree.
\evidence{/Users/maxghenis/PolicyEngine/populace/packages/populace-build/src/populace/build/us\_runtime/reform\_validation.py}
\item \textbf{Tax expenditures vs.\ JCX-48-24 and Treasury OTA.} Five
tax-expenditure reform rows (child tax credit, EITC, child and dependent
care credit, standard deduction, total itemized deductions) carry benchmark
values from \jct{}'s JCX-48-24 Table 1 or Treasury's FY2025 Tax
Expenditures report where a comparable figure exists; two rows
(\texttt{te\_standard\_deduction}, \texttt{te\_itemized\_total}) are marked
not scored because no comparable published figure exists in either source.
\evidence{/Users/maxghenis/PolicyEngine/populace/packages/populace-build/src/populace/build/us/tax\_expenditure\_reforms.json}
One of the five (\texttt{te\_eitc}) is flagged \texttt{in\_sample} (used as
a calibration target rather than an independent check) and must be excluded
from any claim of independent validation.
\item \textbf{Tax expenditure calibration targets vs.\ JCT (2024).} Five
2024 tax-expenditure revenue-loss figures (SALT deduction, medical expense
deduction, charitable deduction, mortgage interest deduction, qualified
business income deduction) enter the calibration target surface directly,
sourced to the Ledger fact stream's JCT family.
\evidence{/Users/maxghenis/PolicyEngine/populace/packages/populace-build/src/populace/build/us/fiscal\_target\_references.json,
\texttt{target\_references} entries with \texttt{family: "jct"}} These are
calibration targets, not held-out validation, and belong in
Section~\ref{sec:calibration} of the model description rather than being
cited as independent evidence here.
\end{itemize}

\todo{Evidence needed: no committed artifact in the surveyed repositories
compares \policyengineus{} to a published \tpc{} estimate. A repository
named \texttt{bottom-50-tax-analysis} compares a \policyengineus{}-derived
income-tax aggregate to \cbo{} baseline receipts projections (2024--2026),
which is a genuine CBO comparison distinct from the JCT comparisons above
and should be surveyed in full before v1 (this outline verified only that
the repository and comparison exist, not the specific numbers).
\evidence{/Users/maxghenis/PolicyEngine/bottom-50-tax-analysis/README.md}
Several repositories use \cbo{} figures only as \emph{inputs}
(uprating targets, behavioral elasticities) rather than as a validation
target and must not be cited as validation evidence: the
\texttt{parameters/calibration/gov/cbo/} parameter tree and
\texttt{cbo\_household\_*} variables in policyengine-us, and
\texttt{cbo-baseline-tracker} and \texttt{deduction-repeal-analysis}, which
compare CBO to CBO or use CBO elasticities as modeling inputs.}

\subsection{\policybench{}: a distinct claim}
\label{sec:validation-policybench}

\policybench{} benchmarks language models' ability to predict tax and
benefit outcomes, using \policyengineus{}'s own output as the reference
value it scores against -- not administrative data.
\evidence{/Users/maxghenis/PolicyEngine/policybench/paper/index.qmd, line
1081: "PolicyBench treats PolicyEngine outputs as benchmark reference
outputs, not as administrative records."} It is therefore evidence of
\policyengineus{}'s use as a trusted reference implementation and of the
model's internal reproducibility infrastructure (versioned, hash-pinned
snapshots of 100 sampled US scenarios against 19 output specs), but it is
\emph{not} independent validation against administrative outturns and must
not be cited as such in this paper.
\evidence{/Users/maxghenis/PolicyEngine/policybench/paper/snapshot/20260501/us\_scenarios.csv;
/Users/maxghenis/PolicyEngine/policybench/policybench/benchmark\_specs.json}
\todo{Decide whether \policybench{} belongs in this section at all, framed
explicitly as "not independent validation," or is better placed in
Section~\ref{sec:model} as evidence of reproducibility tooling. Current
draft keeps it here with an explicit disclaimer so a reader does not
mistake it for administrative validation.}
